\documentclass[10pt,a4paper]{amsart}
\usepackage[margin=19mm]{geometry}
\usepackage{amsmath,amssymb,mathtools}
\usepackage[scr=rsfs,cal=euler]{mathalfa}
\usepackage{xfrac}
\usepackage[unicode,hidelinks,bookmarks=false]{hyperref}
\newcommand{\ie}{i.e.\ }
\newcommand{\cf}{cf.\ }
\newcommand{\resp}{respectively\ }
\newcommand{\Subsection}{\subsection}
\providecommand{\nobreakdash}{\nobreak\hbox{-}}
\setlength{\emergencystretch}{2em}
\multlinegap0pt
\usepackage{bm}
\usepackage{bbm}
\usepackage{dsfont}
\usepackage{xspace}
\usepackage{cancel}
\usepackage{mathtools}
\usepackage{amsthm}
\makeatletter
\@ifundefined{lem}{\newtheorem{lem}{Lemma}}{}
\makeatother
\title[Grassmannian perspectives of classical Lie groups and Cartan]{Grassmannian perspectives of classical Lie groups and Cartan involutions}
\author{Yunxia Chen, Naichung Conan Leung}
\date{Source: arXiv:2602.00602v1 (CC BY 4.0)}
\begin{document}
\maketitle

\noindent\emph{Source and licence.} Grassmannian perspectives of classical Lie groups and Cartan involutions. Version arXiv:2602.00602v1 (CC BY 4.0), distributed under the Creative Commons Attribution 4.0 International licence, as recorded in the arXiv licence metadata for this exact version and shown on the abstract page https://arxiv.org/abs/2602.00602v1. Full rights notice: licenses/arxiv-2602.00602-CC-BY-4.0.txt.\par\smallskip

\noindent\textit{Editorial scope.} Counted unit: the Lem whose source label is \texttt{None} in the article (source0.tex lines 794--797, source label \texttt{None}), delivered together with the article's own complete original proof of it (source0.tex lines 798--890, one \texttt{proof} environment of 93 lines). No mathematics is written, repaired or re-derived: statement and proof are the article's own text line for line. Required context retained uncounted: none. Every definition, display and label of those lines is retained unchanged. Omitted from this package and not redistributed: 1--793, 891--1095. Nothing inside a retained excerpt is omitted or altered: every retained body is delivered exactly as the article's lines are written, so no omission receipt of any kind accompanies this package. Citations: the retained excerpts carry no citation command at all, so no bibliography entry of the article is transcribed and no reference is redistributed. Reference adaptation: none was needed and none was made; every reference inside the delivered unit resolves inside it, so no unresolved reference or citation is produced. Printed slips kept literal: no printed word, symbol, hypothesis or number of the counted statement or of its proof was altered. Layout and class: the article's own class is \texttt{article} with packages including amsfonts, mathrsfs, CJK, amssymb, amsmath, amsthm, graphicx, enumerate, amscd, bbm, verbatim, latexsym, multirow, bigdelim; the wrapper is the lane's pinned \texttt{amsart} editorial wrapper at 10pt on A4, which restates the theorem-like environments the retained text needs and adds the article's own macro definitions that the retained text needs (none), with extra wrapper packages (bm, bbm, dsfont, xspace, cancel, mathtools). The delivered numbering is the unit's own. Assets: no retained span contains a figure environment, a graphics-inclusion command, a PDF-page inclusion, a table or a TikZ picture; no image file is copied into this package, the package declares no asset dependency and no image file is redistributed with it. Licence: version arXiv:2602.00602v1 of this article is distributed under the Creative Commons Attribution 4.0 International licence, as recorded in the arXiv licence metadata for this exact version and shown on the abstract page https://arxiv.org/abs/2602.00602v1. A full rights notice is delivered at licenses/arxiv-2602.00602-CC-BY-4.0.txt. Characters: every retained excerpt is pure ASCII, so the delivered engine, XeLaTeX with its default Latin Modern text font, reports no missing character.
\begin{lem}
For any $L_0, L_1, L_2 \in M_{XY}$, there exists a unique $L' \in M_{XY}$ such that
$\varphi_{L_2}\circ \varphi_{L_0}^{-1}\circ \varphi_{L_1}=\varphi_{L'}$.
\end{lem}

\begin{proof}
We will prove for $\mathbb{F}=\mathbb{R}$, the proof is similar for $\mathbb{F}=\mathbb{C},\mathbb{H}$. Note that this argument only depends on the
relative position of $X$ and $Y$. Since the action of $O(2n)$  on $M=Gr(n,V_1\oplus V_2)$
 is transitive and preserves the relative position of two subspaces, we can assume $X=\left[
\begin{array}{cc}
I \\
O \end{array} \right]$.

From $Y\bigcap X=\{0\}$, we have $Y=\left[
\begin{array}{cc}
A \\
I \end{array} \right]$ with $A$ an $n \times n$ matrix, otherwise we will have $Y=\left[
\begin{array}{cc}
* & *\\
* & 0_{n \times 1}\end{array} \right]$ after elementary column transformation, which is contradict to $Y\bigcap X=\{0\}$.

\bigskip

For any  $L \in M_{XY}$, $L \bigcap X^{\bot}=\{0\}$. Since $X^{\bot}=\left[
\begin{array}{cc}
O \\
I \end{array} \right]$, we have  $L=\left[
\begin{array}{cc}
I \\
B \end{array} \right]$ with $B$ an $n \times n$ matrix. Moreover,  since $Y^{\bot}=\left[
\begin{array}{cc}
I \\
-A^t \end{array} \right]$, we have $L \bigcap Y^{\bot}=\{0\}$ is equivalent to $A^t+B$ is invertible.

\bigskip

Under the above conditions, we have the map $\varphi_{L}$ is given by
\begin{align*}
\varphi_{L} : & ~  ~   ~ X  ~  ~   ~ \longrightarrow  ~  ~   ~ L  ~  ~   ~ ~  ~   ~\longrightarrow ~  ~   ~ ~  ~   ~ Y\\
 ~ ~ & \left(
\begin{array}{cc}
x \\
0 \end{array} \right) \mapsto  \left(
\begin{array}{cc}
x \\
Bx \end{array} \right) \mapsto  \left(
\begin{array}{cc}
A(I+A^tA)^{-1}(A^t+B)x \\
(I+A^tA)^{-1}(A^t+B)x \end{array} \right)
\end{align*}
where the first map from $X$ to $L$ is obvious and the second map from $L$ to $Y$ is given by the projection matrix $P_Y$ of $Y$. By direct computation, we have
\[
P_Y=Y(Y^tY)^{-1}Y^t=\left(
\begin{array}{cc}
A(I+A^tA)^{-1}A^t & A(I+A^tA)^{-1}\\
(I+A^tA)^{-1}A^t & (I+A^tA)^{-1}\end{array} \right).
\]

Now for any  $L_0=\left[
\begin{array}{cc}
I \\
B_0 \end{array} \right], L_1=\left[
\begin{array}{cc}
I \\
B_1 \end{array} \right], L_2=\left[
\begin{array}{cc}
I \\
B_2 \end{array} \right] \in M_{XY}$, we have $A^t+B_i$ is invertible for $i=0,1,2$ and $\varphi_{L_2}\circ \varphi_{L_0}^{-1}\circ \varphi_{L_1}$ is as follows:


\[
X  ~  ~   ~ \longrightarrow  ~  ~   ~ Y  ~  ~   ~ ~  ~   ~\longrightarrow ~  ~   ~ ~  ~   ~ X~  ~   ~ \longrightarrow  ~  ~   ~ Y
\]
\[
\left(
\begin{array}{cc}
x \\
0 \end{array} \right) \mapsto   \left(
\begin{array}{cc}
A(I+A^tA)^{-1}(A^t+B_1)x \\
(I+A^tA)^{-1}(A^t+B_1)x \end{array} \right)
\mapsto   \left(
\begin{array}{cc}
(A^t+B_0)^{-1}(A^t+B_1)x \\
0 \end{array} \right)
\]
\[
 \mapsto   \left(
\begin{array}{cc}
A(I+A^tA)^{-1}(A^t+B_2)(A^t+B_0)^{-1}(A^t+B_1)x \\
(I+A^tA)^{-1}(A^t+B_2)(A^t+B_0)^{-1}(A^t+B_1)x \end{array} \right)
\]
To get $\varphi_{L_2}\circ \varphi_{L_0}^{-1}\circ \varphi_{L_1}=\varphi_{L'}$, we just need to take $L'=\left[
\begin{array}{cc}
I \\
B' \end{array} \right]$ with $B'=(A^t+B_2)(A^t+B_0)^{-1}(A^t+B_1)-A^t$. Note that $A^t+B'=(A^t+B_2)(A^t+B_0)^{-1}(A^t+B_1)$ is invertible,
i.e. $L' \in M_{XY}$. And such $L'$ is unique.
\end{proof}

\end{document}
